%% COMP0222 Coursework 2 -- Group 32
%% Visual and LiDAR SLAM
%% Build: run `git lfs pull` first to fetch result PNGs, then
%%   pdflatex COMP0222_CW2_GRP_32.tex     (run twice so references resolve)
\documentclass[twocolumn,10pt,a4paper]{article}

\usepackage[margin=1.8cm,top=2cm,bottom=2cm]{geometry}
\usepackage{graphicx}
\usepackage{booktabs}
\usepackage{amsmath}
\usepackage{amssymb}
\usepackage{hyperref}
\usepackage{xcolor}
\usepackage{caption}
\usepackage{subcaption}
\usepackage{microtype}
\usepackage{float}
\usepackage{listings}

\captionsetup{font=small, labelfont=bf, skip=3pt}
\setlength{\columnsep}{0.6cm}

\definecolor{codebg}{RGB}{245,245,245}
\lstset{
  basicstyle=\ttfamily\scriptsize,
  backgroundcolor=\color{codebg},
  keywordstyle=\color{blue!70!black},
  commentstyle=\color{green!40!black}\itshape,
  stringstyle=\color{red!60!black},
  breaklines=true,
  breakatwhitespace=true,
  frame=single,
  framesep=2pt,
  rulecolor=\color{gray!40},
  xleftmargin=2pt,
  xrightmargin=2pt,
  aboveskip=3pt,
  belowskip=3pt,
  numbers=none,
  captionpos=b,
  showstringspaces=false,
  columns=fullflexible,
  literate={->}{{$\rightarrow$}}1,
}
\setlength{\parskip}{2pt}

\hypersetup{
  colorlinks=true,
  linkcolor=blue!70!black,
  citecolor=green!50!black,
  urlcolor=blue!70!black
}

\title{\textbf{Visual and LiDAR SLAM}\\
  \large COMP0222 Coursework 2 -- Group 32}
\author{Muhammad Maaz \and Harmish}
\date{April 2026}

\begin{document}
\maketitle

%% ---------------------------------------------------------------------------
\begin{abstract}
\noindent
We run monocular ORB-SLAM2 on KITTI~07 and TUM
\textit{freiburg3\_long\_office}, collect two custom RGB sequences with an
Intel RealSense D455, and build a 2D LiDAR SLAM stack (point-to-plane ICP,
log-odds occupancy, four-gate loop detection, GTSAM pose-graph
optimisation) from RPLidar A2M12 data. Headline numbers: KITTI~07 ATE
4.39\,m (0.6\% of a 695\,m path), TUM ATE
0.032\,m, and LiDAR Basement\_1/Floor7\_Hallway closure errors
below 0.2\,m before pose-graph optimisation. In Q3d, GTSAM reduces
the total pose-graph cost by 8--25$\times$ but slightly worsens the
endpoint closure error (Floor7\_Hallway: 0.166\,m to 0.221\,m),
which we analyse as a loop/odometry weighting issue rather than
presenting as an improvement. Every change to a parameter or a module
has a clear effect that the report explains with the underlying maths.
\end{abstract}

%% ---------------------------------------------------------------------------
\section{Introduction}
\label{sec:intro}
SLAM means: build a map and estimate where the sensor is inside it at the
same time. We study two sensors. An RGB camera gives pixels; we use
ORB-SLAM2 to turn those pixels into a trajectory and a sparse 3D map.
A 2D spinning LiDAR gives range readings; we use point-to-plane ICP and a
factor graph to turn those readings into a 2D trajectory and an occupancy
grid. We organise the work into three questions: ORB-SLAM2 on
benchmark sequences (KITTI~07 and TUM), ORB-SLAM2 on our own
sequences compared with COLMAP, and 2D LiDAR SLAM on our own
sequences with loop closure and pose-graph optimisation.

%% ---------------------------------------------------------------------------
\section{Metrics}
\label{sec:metrics}

\textbf{Absolute Trajectory Error (ATE)} measures global consistency. The estimated trajectory is first aligned to ground truth using a Sim(3) Umeyama transform (rotation, translation, and scale), which is necessary for monocular SLAM where the absolute scale is unobservable. ATE RMSE is then the average per-pose translation error after this alignment. A low ATE means the overall shape of the trajectory matches ground truth.

\textbf{Relative Pose Error (RPE)} measures local, per-frame accuracy. It computes how much the estimated motion between two consecutive frames deviates from the true motion, independent of any global alignment. A low RPE means individual frame-to-frame steps are accurate, even if the trajectory drifts globally.

\textbf{LiDAR odometry} uses scan-to-grid matching: each new scan is matched against the current occupancy map by searching for the pose correction that maximises the number of scan points landing on occupied cells. The map is updated with log-odds: each scan hit increases a cell's occupancy score by a fixed amount, each beam pass-through decreases it. Huber-weighted point-to-plane ICP provides the covariance estimate for each odometry step, which is passed to the pose-graph as an edge weight.

\textbf{Pose-graph optimisation} collects all odometry edges and detected loop closure edges into a factor graph and minimises the total weighted inconsistency using Levenberg--Marquardt. Loop edges are only trusted in proportion to their ICP alignment quality, so weak loop matches have little influence on the result.

%% ===========================================================================
\section{Q1: ORB-SLAM2 on Benchmark Datasets}
\label{sec:q1}

\subsection*{Datasets and Protocol}

KITTI~07 is a stereo driving sequence recorded in Karlsruhe, Germany; we run ORB-SLAM2 on the left camera only in monocular mode. The path is 694.7\,m over 1101~frames at 10\,fps with one prominent loop near the end, and ground truth comes from a GNSS/IMU rig at centimetre accuracy. TUM \textit{freiburg3\_long\_office\_household} is an indoor office scan captured handheld with an RGB-D camera, of which we use only the RGB stream. The path is 22.2\,m over 2557~frames at 30\,fps, with sub-millimetre motion-capture ground truth.

All trajectories are evaluated using EVO with Sim(3) Umeyama alignment to account for the inherent scale ambiguity of monocular SLAM. For each ablation (Q1b--Q1d) we change exactly one parameter or code path and re-run from scratch.

\subsection*{Q1(a) Baseline Evaluation}

\begin{table}[H]
\centering\small
\caption{Baseline ATE and RPE for both sequences.}
\label{tab:q1a}
\resizebox{\columnwidth}{!}{%
\begin{tabular}{lcccc}
\toprule
Sequence & Poses & ATE RMSE & ATE \% path & RPE RMSE \\
\midrule
KITTI~07              & 1096 & 4.39\,m  & 0.63\,\% & 0.675\,m \\
TUM fr3\_long\_office & 2557 & 0.032\,m & 0.14\,\% & 0.006\,m \\
\bottomrule
\end{tabular}%
}
\end{table}

\begin{figure}[H]
  \centering
  \includegraphics[width=\columnwidth]{plots/q1a_baseline.png}
  \caption{Baseline trajectories and per-frame ATE for KITTI~07 (left) and TUM (right).}
  \label{fig:q1a}
\end{figure}

\begin{figure}[H]
  \centering
  \includegraphics[width=0.49\columnwidth]{plots/q1a_kitti_ape_map.png}
  \hfill
  \includegraphics[width=0.49\columnwidth]{plots/q1a_tum_ape_map.png}
  \caption{EVO APE colour maps. Error colour scale from low (blue) to high (red). KITTI shows increasing error towards the end of the loop; TUM errors remain low throughout.}
  \label{fig:q1a_evo}
\end{figure}

The baseline results reflect the structural differences between the two sequences. On KITTI~07, the ATE of 4.39\,m represents 0.63\,\% of the 695\,m drive, which is consistent with published monocular ORB-SLAM2 performance on this sequence. The per-frame ATE plot shows drift accumulating gradually through the drive, with a sharp correction at the loop closure event near frame 1000. The RPE of 0.675\,m/frame indicates that the majority of absolute error originates from accumulated heading drift rather than from noisy individual frame estimates.

On TUM, the ATE of 0.032\,m (0.14\,\% of a 22\,m path) is near-optimal for monocular SLAM. The slow handheld motion produces small inter-frame baselines well-suited to triangulation, the richly textured office environment provides abundant distinctive ORB keypoints, and the confined spatial scale limits the distance over which drift can grow. The contrast between the two sequences demonstrates that environment scale and motion speed are the primary determinants of monocular SLAM accuracy, independent of the algorithm configuration.

\subsection*{Q1(b) Number of ORB Features}

The number of features extracted per frame is controlled by a single parameter in the YAML configuration file:

\begin{lstlisting}[language=bash]
# e.g. KITTI04-12_custom_f2000.yaml
ORBextractor.nFeatures: 2000
\end{lstlisting}

We tested three feature levels per sequence. For KITTI: 1500, 2000, and 2500 features; all three remain above the monocular initialisation threshold and track the full sequence. For TUM: 500, 800, and 1000 features; 250 features fails initialisation entirely.

\begin{table}[H]
\centering\small
\caption{Feature count effect on ATE and tracking rate for KITTI~07 and TUM.}
\label{tab:q1b}
\resizebox{\columnwidth}{!}{%
\begin{tabular}{lcccc|lcccc}
\toprule
\multicolumn{5}{c|}{KITTI 07} & \multicolumn{5}{c}{TUM fr3\_long\_office} \\
Features & ATE (m) & RPE (m) & Poses & Track. & Features & ATE (m) & RPE (m) & Poses & Track. \\
\midrule
2500 & 17.55 & 0.150 & 1097 & 100\,\% & 1000\,(base) & 0.032 & 0.006 & 2557 & 100\,\% \\
2000 & 18.47 & 0.155 & 1096 & 100\,\% & 800          & 0.109 & 0.006 & 2556 & 100\,\% \\
1500 & 18.74 & 0.162 & 1096 & 100\,\% & 500          & 0.015\,$^\ast$ & 0.007 & 1152 & 45\,\% \\
\bottomrule
\end{tabular}%
}
\\[3pt]
\footnotesize $^\ast$ feat=500 ATE appears lower than baseline because 55\,\% of frames are lost before drift accumulates --- a measurement artefact, not a genuine improvement.
\end{table}

\begin{figure}[H]
  \centering
  \includegraphics[width=\columnwidth]{plots/q1b_kitti_features.png}
  \caption{KITTI07 trajectories at three feature levels. All three track the full sequence; RPE decreases monotonically as features increase.}
  \label{fig:q1b_kitti}
\end{figure}

\begin{figure}[H]
  \centering
  \includegraphics[width=\columnwidth]{plots/q1b_tum_features.png}
  \caption{TUM trajectories at three feature levels plus ATE trend. feat=250 fails initialisation entirely.}
  \label{fig:q1b_tum}
\end{figure}

\textbf{KITTI07.} All three configurations track the complete 1096--1097 pose sequence. Per-frame accuracy improves monotonically with feature count: RPE RMSE falls from 0.162\,m at feat=1500 to 0.150\,m at feat=2500, a 7\,\% improvement. ATE is approximately 18\,m for all three runs because loop closure did not fire consistently during these trials --- a result in line with our Q1(d) no-loop baseline of 18.02\,m. This reveals a fundamental decoupling: \emph{feature count governs per-frame tracking quality (RPE), while global trajectory consistency (ATE) is determined by whether loop closure fires}. More features improve local estimation but cannot correct accumulated drift in the absence of a closure event.

\textbf{TUM.} Reducing features to 800 degrades ATE by 3.4$\times$ (0.109\,m vs 0.032\,m) while maintaining full tracking. Fewer descriptors produce more ambiguous nearest-neighbour matches during pose optimisation, allowing heading bias to accumulate over the 22\,m path. At feat=500, tracking is lost for 55\,\% of frames: the system cannot maintain the minimum feature density required for reliable pose estimation across the indoor office environment. At feat=250, initialisation fails entirely.

\subsection*{Q1(c) Outlier Rejection}

ORB-SLAM2's pose optimisation (\texttt{Optimizer::PoseOptimization}) runs four rounds of Levenberg--Marquardt bundle adjustment. After each round, feature correspondences whose reprojection residual exceeds a chi-squared threshold are flagged as outliers and excluded from the next round. We disabled this gate in \texttt{src/Optimizer.cc} by setting a single boolean:

\begin{lstlisting}[language=C++]
const bool kDisableOutlierRejection = true;
...
if (!kDisableOutlierRejection && chi2 > chi2Mono[it]) {
    pFrame->mvbOutlier[idx] = true;
    e->setLevel(1);
}
\end{lstlisting}

With the constant set to \texttt{true}, the condition is never entered and every feature match --- including false ones --- is trusted by the LM solver.

\begin{table}[H]
\centering\small
\caption{Effect of disabling outlier rejection on both sequences.}
\label{tab:q1c}
\resizebox{\columnwidth}{!}{%
\begin{tabular}{lcccc}
\toprule
Configuration & ATE RMSE & RPE RMSE & Poses tracked & ATE change \\
\midrule
KITTI baseline   & 4.39\,m  & 0.675\,m & 1096 & --- \\
KITTI no outlier & 5.77\,m  & 0.825\,m & 533  & $+31\,\%$ \\
TUM baseline     & 0.032\,m & 0.006\,m & 2557 & --- \\
TUM no outlier   & 0.470\,m & 0.053\,m & 1223 & $+1380\,\%$ \\
\bottomrule
\end{tabular}%
}
\end{table}

\begin{figure}[H]
  \centering
  \includegraphics[width=\columnwidth]{plots/q1c_outlier.png}
  \caption{Outlier rejection ablation. Red: no outlier rejection. Blue: baseline. Both sequences lose approximately half their tracked poses.}
  \label{fig:q1c}
\end{figure}

Disabling outlier rejection had a substantially larger impact on TUM than on KITTI, which reflects the fundamental difference in scene content. The KITTI outdoor environment is predominantly static --- road surfaces, buildings, and open sky --- providing very few false matches for the pose solver. On TUM, glass surfaces, reflective floors, and intermittent moving people regularly violate the rigid-world assumption underlying bundle adjustment. Without the outlier gate, every erroneous correspondence is weighted equally in the normal equations, corrupting the pose estimate and compounding error frame by frame.

The impact is quantified in both metrics. ATE on TUM increases by 1380\,\% (0.032 to 0.470\,m) compared to only 31\,\% on KITTI. RPE on TUM rises ninefold (0.006 to 0.053\,m), confirming that per-step motion is directly corrupted rather than just global consistency. Critically, both sequences also lose approximately half their tracked poses: corrupted pose estimates push the system outside its relocalisation basin, causing hard resets. The no-outlier mode is therefore both less accurate and substantially less robust, with indoor scenes suffering far more severely due to the prevalence of dynamic and specular elements.

\subsection*{Q1(d) Loop Closure}

ORB-SLAM2's loop closure thread (\texttt{LoopClosing}) continuously queries a DBoW2 bag-of-words database for keyframes that share visual similarity with the current keyframe. When a geometric consistency check passes, it computes a 7-DoF Sim(3) correction and redistributes accumulated drift across the full pose chain. We disabled this thread by forcing \texttt{DetectLoop} to return immediately at entry:

\begin{lstlisting}[language=C++]
// LoopClosing.cc (Q1d patch)
const bool kDisableLoopClosure = true;
if (kDisableLoopClosure) {
    mpKeyFrameDB->add(mpCurrentKF);
    mpCurrentKF->SetErase();
    return false;
}
\end{lstlisting}

Tracking and local bundle adjustment continue to run normally; only the global loop correction is suppressed.

\begin{table}[H]
\centering\small
\caption{Effect of disabling loop closure on both sequences.}
\label{tab:q1d}
\resizebox{\columnwidth}{!}{%
\begin{tabular}{lcccc}
\toprule
Configuration & ATE RMSE & RPE RMSE & Poses tracked & ATE change \\
\midrule
KITTI baseline  & 4.39\,m  & 0.675\,m & 1096 & --- \\
KITTI no loop   & 18.02\,m & 0.154\,m & 1096 & $+310\,\%$ \\
TUM baseline    & 0.032\,m & 0.006\,m & 2557 & --- \\
TUM no loop     & 0.049\,m & 0.006\,m & 2557 & $+54\,\%$ \\
\bottomrule
\end{tabular}%
}
\end{table}

\begin{figure}[H]
  \centering
  \includegraphics[width=\columnwidth]{plots/q1d_loop.png}
  \caption{Loop closure ablation. Orange: no loop closure. Blue: baseline. KITTI shows severe drift in the second half of the sequence; TUM shows only modest degradation.}
  \label{fig:q1d}
\end{figure}

Disabling loop closure reveals a clean dissociation between local and global error. On KITTI~07, ATE jumps from 4.39\,m to 18.02\,m ($+310\,\%$) as the heading drift accumulated over 695\,m is never corrected. The trajectory fails to close its loop and diverges progressively from ground truth throughout the second half of the drive. On TUM, the impact is comparatively small (ATE rises by 54\,\%, from 0.032 to 0.049\,m) because there is far less drift to correct over a 22\,m indoor path.

The most revealing observation is the behaviour of RPE. On KITTI, RPE \emph{decreases} from 0.675\,m to 0.154\,m without loop closure. This occurs because loop correction must rewrite historical pose estimates to satisfy the closure constraint, slightly disrupting the smoothness of local frame-to-frame transitions. Without loop closure, local tracking proceeds undisturbed, improving RPE while simultaneously allowing ATE to deteriorate. This confirms that ATE and RPE measure fundamentally different properties: ATE reflects accumulated global drift, while RPE reflects per-frame matching quality.

Taken together, Q1(c) and Q1(d) produce opposite failure signatures. Disabling outlier rejection inflates RPE without necessarily destroying global consistency; disabling loop closure inflates ATE while leaving RPE unchanged or improved. Each component of ORB-SLAM2 therefore addresses a distinct failure mode, and both are required for reliable operation across diverse environments.

\subsection*{Q1 Summary}

\begin{figure}[H]
  \centering
  \includegraphics[width=\columnwidth]{plots/q1_rpe.png}
  \caption{RPE across all Q1 configurations. Outlier removal produces the largest RPE excursions; loop removal sits at or below baseline because per-frame matching is unaffected.}
  \label{fig:q1_rpe}
\end{figure}

The four ablations characterise ORB-SLAM2's architecture at the component level. The feature budget sets the quality ceiling for per-frame tracking. The outlier gate prevents corrupted matches from propagating into the pose estimate at tracking time. Loop closure is the only mechanism that corrects globally accumulated drift. These components are independent: removing any one degrades a distinct aspect of performance, and no single metric captures all three failure modes simultaneously.

%% ===========================================================================
\section{Q2: ORB-SLAM2 on Custom Sequences}
\label{sec:q2}

We collect two sequences, one indoor and one outdoor, both with at
least 500 tracked poses after ORB-SLAM2 initialisation:
\textbf{Basement\_1} (indoor, 869 poses) and \textbf{Outdoor\_1}
(outdoor, 983 poses).

\subsection{Q2a: Data collection and calibration}
\label{sec:q2a}

The recording rig is an Intel RealSense D455 colour stream at
$848 \times 480$ and 30\,fps. Auto-exposure was locked after a
two-second warm-up so the intrinsics did not drift during recording.
The same colour stream feeds both ORB-SLAM2 monocular and COLMAP.

\textbf{Basement\_1} is the lower-ground-floor Lift D room at Marshgate
--- a compact space containing lift doors and fire-exit doors, with
concrete walls on all four sides and two doorways. The enclosed
geometry provides COLMAP with stable keypoints on every wall surface
for robust triangulation. ORB-SLAM2 produced 869 tracked poses over a
closed loop that returns to the start.

\textbf{Outdoor\_1} is the back-entrance cycle-parking area of the
Marshgate building, approached from the service exit. The right side
is a flat brick building fa\c{c}ade; the left side opens onto a row
of trees and a paved path. ORB-SLAM2 produced 983 tracked poses along
a single rectangular loop around the courtyard.

The capture strategy was chosen to maximise the overlap between
consecutive frames so that both COLMAP and ORB-SLAM2 can match
keypoints reliably. We walked at 0.3 to 0.5\,m/s, which at 30\,fps
gives 1 to 1.5\,cm of camera translation per frame, well below the
focal-length-scaled disparity needed for ORB to match the same
keypoint across consecutive frames. We recorded a continuous video
stream rather than still shots, because a still-shot sequence would
force the initialiser to choose two frames with a wide baseline and
that breaks ORB matching. At each corner we paused for one second
and panned slowly through 90 degrees, which gives COLMAP the
multi-view diversity its bundle adjustment needs to converge. The
ORB-SLAM2 viewer was monitored live throughout the recording, and we
slowed down whenever the tracked-feature count fell below about 30.

Calibration uses COLMAP in two ways. The first is per-sequence
refinement with the OPENCV 8-parameter model
$(f_x, f_y, c_x, c_y, k_1, k_2, p_1, p_2)$, where each sequence is
run through COLMAP, the resulting \texttt{cameras.txt} is converted
to an ORB-SLAM2 YAML with \texttt{q2\_calibration\_report.py}, and
ORB-SLAM2 is re-run with the refined intrinsics; calibrated
trajectories are saved as
\texttt{\{seq\}\_trajectory\_colmap\_intrinsics.txt} and are used for
all Q2b comparisons as required by the brief. The second calibration
approach is a genuine no-prior run on Basement\_1 with the
SIMPLE\_PINHOLE model ($f, c_x, c_y$, no distortion) and no factory
prior at all, confirming the intrinsics can be recovered from the
images alone.

\begin{table}[H]
\centering\scriptsize
\caption{Q2a COLMAP self-calibration vs.\ factory D455 intrinsics.}
\label{tab:q2a_calib}
\resizebox{\columnwidth}{!}{%
\begin{tabular}{lccc}
\toprule
Run & $f_x$ (px) & $f_y$ (px) & $\Delta f_x$ \\
\midrule
Factory D455                          & 426.68 & 426.10 & n/a \\
Basement\_1 scratch (SIMPLE\_PINHOLE) & 415.40 & 415.40 & $-2.6\,\%$ \\
Basement\_1 (OPENCV)                  & 430.64 & 445.84 & $+0.9\,\%$ \\
Outdoor\_1  (OPENCV)                  & 427.91 & 437.20 & $+0.3\,\%$ \\
\bottomrule
\end{tabular}%
}
\end{table}

Both OPENCV runs agree with the factory $f_x$ within 1\,\%,
confirming the D455 factory calibration is accurate for $f_x$.
However, $f_y$ is inflated by 4.6\,\% and 2.6\,\% above factory
--- a known COLMAP artefact where the bundle adjuster trades $f_y$
against the radial distortion coefficients to minimise reprojection
error when scene depth variation is limited; the physical $f_y$ is
not recovered cleanly. The scratch SIMPLE\_PINHOLE result
($f = 415.4$\,px) is 2.6\,\% below factory because SIMPLE\_PINHOLE
forces $f_x = f_y$ while the D455 is mildly asymmetric; this is a
model constraint, not a calibration error. The consequences of the
$f_y$ inflation on ORB-SLAM2 are analysed in Q2b.

Figure~\ref{fig:q2b_colmap} (next subsection) shows both ORB-SLAM2
trajectories coloured by elapsed time.

\subsection{Q2b: Results}
\label{sec:q2b}

No external ground truth exists for either custom sequence, so we
cannot compute a true ATE. What we report is inter-method
disagreement: how far the ORB-SLAM2 trajectory is from the COLMAP
trajectory after a Sim(3) Umeyama alignment with scale (Eq.~1, 2).
A small number means the two independent monocular pipelines agree on
the trajectory shape, but it does not mean either trajectory is
correct in metric units.

For each sequence we load the COLMAP keyframe poses (from
\texttt{images.txt}) and the ORB-SLAM2 trajectory (TUM format). The
real RGB timestamp is re-attached to each COLMAP keyframe by looking
its image filename up in \texttt{rgb.txt}. EVO's
\texttt{sync.associate\_trajectories} then pairs COLMAP keyframes
with ORB-SLAM2 frames within 30\,ms (relaxed to 250\,ms if the
strict tolerance returns no matches), and an EVO Umeyama alignment
with scale puts both trajectories in the same frame. We report APE
in three pose relations: translation only, rotation only (degrees),
and full SE(3).

\begin{table}[H]
\centering\scriptsize
\caption{Q2b inter-method disagreement (COLMAP vs.\ ORB-SLAM2 with
COLMAP-calibrated intrinsics, EVO Sim(3) alignment). Matched is the
number of timestamp-aligned pose pairs after MAD outlier filtering.}
\label{tab:q2b}
\resizebox{\columnwidth}{!}{%
\begin{tabular}{lcccc}
\toprule
Sequence & ORB poses & Matched & Trans.\ RMSE & Rot.\ RMSE \\
\midrule
\textbf{Basement\_1} (indoor)  & 868 & 44 & \textbf{0.046\,a.u.} & $3.8^{\circ}$ \\
\textbf{Outdoor\_1}  (outdoor) & 584 & 54 & \textbf{1.488\,a.u.} & $169.2^{\circ}$ \\
\bottomrule
\end{tabular}%
}
\end{table}

RMSE values are in arbitrary units (a.u.) because both monocular
methods recover trajectories up to an unknown scale factor; the
Sim(3) alignment absorbs scale but cannot assign metric units.

\begin{figure}[H]
  \centering
  \includegraphics[width=\columnwidth]{plots/q2b_colmap_trajectories.png}
  \caption{Q2b COLMAP SfM trajectories coloured by elapsed time.}
  \label{fig:q2b_colmap}
\end{figure}

\begin{figure}[H]
  \centering
  \includegraphics[width=\columnwidth]{plots/q2b_colmap_vs_orbslam.png}
  \caption{Q2b COLMAP (pseudo-GT, left) vs.\ ORB-SLAM2 (EVO-aligned,
  right) for each sequence. Gray ghost shows the full COLMAP
  reconstruction; coloured trajectory is the timestamp-matched subset.}
  \label{fig:q2b}
\end{figure}

Figure~\ref{fig:q2b_colmap} shows the raw COLMAP SfM trajectories
coloured by elapsed recording time. Basement\_1 traces a short
closed loop around the Lift D room (283 keyframes); Outdoor\_1 traces
a larger rectangular courtyard path (360 keyframes).

Figure~\ref{fig:q2b} places COLMAP and the EVO-aligned ORB-SLAM2
trajectory side by side. The gray ghost is the full COLMAP
reconstruction; the coloured gradient shows the timestamp-matched
subset used for EVO comparison.

\textbf{Basement\_1 analysis.}
COLMAP selected 283 keyframes from the 868-pose calibrated
ORB-SLAM2 log ($\approx$33\,\% keyframe rate). Of those 283, 44
(16\,\%) produced a timestamp match within 30\,ms after MAD outlier
rejection removed 8 pairs. The low match-rate reflects the compact
room geometry: the camera revisits the same viewpoints repeatedly
and ORB-SLAM2 suppresses duplicate keyframe insertion, leaving gaps
in the timestamp-aligned subset. The 44 clean pairs yield
0.046\,a.u.\ translation RMSE and $3.8^{\circ}$ rotation
disagreement --- both COLMAP and ORB-SLAM2 converge to nearly
identical loop shapes in this well-textured indoor space, confirming
that the small COLMAP $f_y$ inflation ($+4.6\,\%$) does not
meaningfully degrade indoor tracking.

\textbf{Outdoor\_1 analysis.}
COLMAP selected 360 keyframes from the 584-pose calibrated
ORB-SLAM2 log (37\,\% keyframe rate). Timestamp association
returned 54 matched pairs after MAD outlier rejection removed 4
pairs.

The $169.2^{\circ}$ rotation disagreement is caused by a
monocular initialisation flip triggered by the COLMAP $f_y$
inflation. The mechanism is as follows. ORB-SLAM2 initialises by
triangulating the first map points from two keyframes. Each pixel
is back-projected into a 3D ray using $f_y$; with
$f_y = 437.2$\,px instead of the true $\approx 426.1$\,px, every
vertical pixel shift is interpreted as corresponding to a shallower
depth than it actually is. This compresses the estimated vertical
parallax and biases the perceived scene structure. In an open
outdoor scene with limited close-range geometry the initialiser has
no other constraint to override this bias, so it resolves the
monocular depth ambiguity into the mirror-image solution ---
rotating the entire coordinate frame by $\approx 180^{\circ}$.
Once the first map is initialised with a flipped frame, all
subsequent tracking is consistent within that flipped frame, so
ORB-SLAM2 completes 584 poses without further failure; the flip
only becomes visible when the trajectory is compared against
COLMAP's independently initialised frame.

\begin{table}[H]
\centering\scriptsize
\caption{Outdoor\_1 diagnostic: calibrated vs.\ factory intrinsics.}
\label{tab:q2b_diag}
\resizebox{\columnwidth}{!}{%
\begin{tabular}{lccc}
\toprule
Intrinsics & $f_y$ (px) & Matched & Rot.\ RMSE \\
\midrule
COLMAP OPENCV (calibrated) & 437.2 & 54  & $169.2^{\circ}$ \\
Factory D455               & 426.1 & 89  & $\mathbf{0.9^{\circ}}$ \\
\bottomrule
\end{tabular}%
}
\end{table}

Table~\ref{tab:q2b_diag} isolates the cause: switching to factory
$f_y$ on the same sequence and same COLMAP reference drops the
rotation disagreement from $169.2^{\circ}$ to $0.9^{\circ}$
and increases matched pairs from 54 to 89. The scene geometry is
well-conditioned; the failure is entirely attributable to the
$f_y$ error in the COLMAP calibration.

The root cause of the $f_y$ inflation is the choice of calibration
data. COLMAP was run on the full walking sequence, where all views
are roughly on the same plane and depth variation is limited. The OPENCV
bundle adjuster then trades $f_y$ against the distortion
coefficients to minimise reprojection error, inflating $f_y$ beyond
its physical value. The Lab~07 guidelines recommend using a
dedicated subset of frames where the camera translates while facing
a planar wall, which provides the focused multi-view geometry needed
to constrain $f_y$ independently of distortion. We also tested that
Lab~07-style subset calibration on our recordings; however, in this
dataset it produced less stable intrinsics and worse downstream
ORB-SLAM2 agreement than the full-sequence calibration. So while the
guideline is correct in principle, our empirical result here is that
the wall-subset variant was noisier and did not improve the Outdoor\_1
failure mode.

\begin{figure}[H]
  \centering
  \includegraphics[width=\columnwidth]{plots/q2b_pointcloud_outdoor_1.png}
  \caption{Q2b COLMAP sparse map, Outdoor\_1 (17{,}806 points).}
  \label{fig:q2b_pcd}
\end{figure}

\subsection{Q2c: Video}
\label{sec:q2c}

The submitted video file \texttt{COMP0222\_CW2\_GRP\_32.mp4} shows,
side by side for both sequences, the live RealSense camera feed and
the ORB-SLAM2 Pangolin viewer reconstructing the scene in real time.
The trajectory draws itself pose by pose as the recording plays, so
the viewer can see the moment of initialisation, every keyframe
insertion, and the final shape of each loop. The Q2 portion is the
first 30 seconds; the Q3 LiDAR portion follows.

Three limitations apply to the Q2 numbers. Monocular SLAM has no
metric scale, so all distances are in arbitrary units. No external
ground truth exists, so we report inter-method agreement rather than
accuracy. The Outdoor\_1 calibrated YAML triggered a $169^{\circ}$
heading flip due to COLMAP OPENCV $f_y$ inflation; running with
factory intrinsics as a diagnostic confirms $0.9^{\circ}$ agreement,
so the failure is attributable to the calibration artefact rather
than scene or algorithm quality.

%% ===========================================================================
\section{Q3: 2D LiDAR SLAM}
\label{sec:q3}

We process three sequences end to end: two indoor (Basement\_1 and
Floor7\_Hallway, the second a long Marshgate corridor that serves as
the large-area indoor case) and one outdoor (Outdoor\_1, an open
courtyard). The same physical environments were used in Q2.

\subsection{Q3a: Data collection}
\label{sec:q3a}

The sensor is an RPLidar A2M12, a single-line spinning 2D LiDAR with a
$360^{\circ}$ scan, 12\,m rated maximum range, up to 16k samples
per second, and a scan rate of roughly 10--12\,Hz in our logs. Each return arrives as
a triple $(q, \alpha, d)$ where $q$ is the return quality (0 to 63),
$\alpha$ is the bearing angle in degrees, and $d$ is the range in
millimetres. We save scans to a JSON-Lines file, one scan per line.

Every scan is filtered before processing. We discard returns with
quality $q < 5$ because low-quality returns come from multipath and
weak surfaces such as glass or polished floors. We mask the
operator's body by dropping
$\alpha \in [135^{\circ}, 225^{\circ}]$, because the operator's legs
sit in that arc during walking. A range gate
$10\,\text{mm} \le d \le \texttt{max\_range}$ removes invalid short
returns and clipped maximum-range returns.

Three environments were used.
\textbf{Basement\_1} is the entrance lobby of the Marshgate building.
Entering through the front door, a staircase door is immediately on
the left and two lifts are on the right; the sequence was collected
in this lobby space, which forms a single rectangular area
approximately 7\,m $\times$ 3\,m. The two open doorways (staircase
and main entrance, visible as free-space fans in the occupancy grid,
Figure~\ref{fig:q3a_grids}) provide two escape directions for the
LiDAR beam, while the solid walls on the remaining sides give
strong 360$^{\circ}$ wall coverage, making this the best-conditioned
sequence for ICP.
\textbf{Floor7\_Hallway} is a wide straight corridor on the seventh
floor of Marshgate, serving as the required large-area indoor scene.
The corridor extends well beyond 15\,m in both directions; multiple
side passages and alcoves branch off the main axis (visible as
long free-space arms in the occupancy grid,
Figure~\ref{fig:q3a_grids}~right). Because the corridor has no room
boundary to loop around, the two-loop route was walked as a compact
square circuit (${\approx}3\,\text{m} \times 2.5\,\text{m}$) at a
widened section so that both loops are geometrically closed while the
LiDAR still observes the full corridor width in every direction.
\textbf{Outdoor\_1} is an open-air courtyard at the rear exit of the
Marshgate building, with the building fa\c{c}ade on one side and
trees and open grass on the other. This mixed geometry---one hard
vertical surface and one soft irregular boundary---means that the
sensor only receives strong returns from one side; the normal
equations for point-to-plane ICP become nearly rank-deficient and
the trajectory freezes (failure analysis below).

For every sequence we marked the floor with chalk, started exactly
on the marker, walked the route at 0.4 to 0.8\,m/s, completed two
clockwise loops, and returned to the marker. The marker was removed
at the end so that the recording itself contains no visual cue of
the start point. Walking pace was chosen so that consecutive scans
overlap by 90\,\% or more, which keeps ICP inside its convergence
basin.

The pipeline runs as follows for each scan. Load the scan and apply
the filters above. Convert to Cartesian with $x = d\cos\alpha$,
$y = d\sin\alpha$. Voxel-downsample at 0.05\,m using the cell
centroid. Estimate per-point surface normals by PCA over $k = 5$
nearest neighbours (smallest eigenvector of the local covariance).
Run point-to-plane ICP  against a rolling local
map made of the last 20 keyframes, with Huber loss $\delta = 0.3$\,m
and IRLS reweighting based on the median absolute deviation. Save the
ICP Hessian $H = J^\top W J$ as the inverse covariance on
$(\theta, t_x, t_y)$ for the pose graph in Q3d. Reject the step if
$\|\Delta t\| > 0.5$\,m or $|\Delta\theta| > 25^{\circ}$ between
consecutive scans, since walking pace cannot exceed these values.
Update the log-odds occupancy grid by Bresenham ray casting from the
sensor pose to each return . Insert a new
keyframe when $\|\Delta t\| > 0.2$\,m or $|\Delta\theta| > 0.2$\,rad
relative to the last keyframe.

The main configuration used for Q3a, Q3c and Q3d is
$\texttt{max\_range} = 12000$\,mm, $\texttt{angular\_step} = 1$
(all beams), $\texttt{voxel} = 0.05$\,m, $\texttt{scan\_skip} = 1$
(all scans); Q3b explores why each value matters. We define closure error as the straight-line distance between the first and last 2D pose.
Two-loop verification uses proximity: a new lap is recorded when the
trajectory enters a 1.5\,m disc around the origin after travelling at
least 5\,m from the last lap boundary, and cumulative heading
integrated from ICP's $\Delta\theta$ is expected near
$\pm 720^{\circ}$ for two clean clockwise loops.

\begin{table}[H]
\centering\scriptsize
\caption{Q3a sequence summary and closure error from the main
configuration.}
\label{tab:q3a}
\resizebox{\columnwidth}{!}{%
\begin{tabular}{lcccc}
\toprule
Sequence & Scans & Keyframes & Path (m) & Closure (m) \\
\midrule
Basement\_1     & 1939 & 172 & 35.0 & 0.17 \\
Floor7\_Hallway & 1304 & 135 & 26.8 & 0.16 \\
Outdoor\_1      & 1462 &  52 & 19.9 & 8.17 \\
\bottomrule
\end{tabular}%
}
\end{table}

\begin{figure}[H]
  \centering
  \includegraphics[width=\columnwidth]{plots/q3a_two_loop_basement_1.png}
  \caption{Q3a Basement\_1 two-loop check.}
  \label{fig:q3a_basement}
\end{figure}

\begin{figure}[H]
  \centering
  \includegraphics[width=0.49\columnwidth]{plots/q3a_two_loop_floor7_hallway.png}
  \hfill
  \includegraphics[width=0.49\columnwidth]{plots/q3a_two_loop_outdoor_1.png}
  \caption{Q3a Floor7\_Hallway and Outdoor\_1 two-loop checks.}
  \label{fig:q3a_others}
\end{figure}

Basement\_1 and Floor7\_Hallway both close within 20\,cm of the
start, with cumulative heading near $-720^{\circ}$ for two clockwise
loops, exactly what the protocol asks for. Outdoor\_1 is the failure
case: the right panel of Figure~\ref{fig:q3a_others} shows the
trajectory stuck at a single point from scan 400 onward. The open
courtyard provides too few stable surfaces for ICP, so
correspondences collapse, the Hessian becomes singular, and the
divergence gate (0.5\,m, 25 degree per-step cap) repeatedly rejects
the step; the pose stops updating but the scan counter continues. We
keep this sequence as the outdoor case and treat the failure mode as
its own finding rather than hide it.

\begin{figure}[H]
  \centering
  \includegraphics[width=0.49\columnwidth]{plots/occupancy_basement_1.png}
  \hfill
  \includegraphics[width=0.49\columnwidth]{plots/occupancy_floor7_hallway.png}
  \caption{Q3a occupancy grids from the main configuration. Left:
    Basement\_1 --- a single rectangular room (${\approx}7\,\text{m}
    \times 3\,\text{m}$) with free-space fans visible through two open
    doorways. Right: Floor7\_Hallway --- a wide corridor extending
    beyond the 10\,m plot boundary in both directions; side passages
    branch off the main axis. The compact square trajectory (red) sits
    at a widened section of the corridor.}
  \label{fig:q3a_grids}
\end{figure}

\subsection{Q3b: Laser odometry parameter analysis}
\label{sec:q3b}

We vary four parameters one at a time on all three sequences. Each
table reports closure error (straight-line distance between first and last pose) per sequence and
the figure under each parameter is the per-parameter occupancy grid
for Basement\_1, the cleanest sequence, which shows each geometric
mechanism most clearly.

\subsubsection*{Maximum range}

\begin{table}[H]
\centering\scriptsize
\caption{Q3b maximum-range sweep. Closure error in metres.}
\label{tab:q3b_range}
\resizebox{\columnwidth}{!}{%
\begin{tabular}{lccc}
\toprule
Setting           & Basement\_1 & Floor7\_Hallway & Outdoor\_1 \\
\midrule
2000\,mm          & 2.38 & 2.02 & 2.81 \\
12000\,mm (sensor max) & \textbf{0.19} & \textbf{0.16} & 8.43 \\
\bottomrule
\end{tabular}%
}
\end{table}

\begin{figure}[H]
  \centering
  \includegraphics[width=\columnwidth]{plots/q3b_grids_basement_1_max_range.png}
  \caption{Basement\_1 grids at 2000 and 12000\,mm.}
  \label{fig:q3b_range}
\end{figure}

ICP  needs surface normals on both sides of the
motion direction to lock translation. With only the near wall in
range, all returns lie on a single line, so the point-to-plane normal
equations become rank-deficient along that line and the trajectory
slides; the $12\times$ improvement on Basement\_1 (2.38 to 0.19\,m
closure) follows directly from this geometric constraint. The
opposite happens outdoors: at 12000\,mm the returns include far,
weakly supported points beyond ICP's recovery basin, so wrong
correspondences pull the registration. The 2000\,mm setting on
Outdoor\_1 gives a smaller closure error (2.81\,m) than 12000\,mm
(8.43\,m) because the shorter range removes the noisy distant
returns. The lesson is that maximum range should be set to roughly
the largest stable surface in view, not always to the sensor maximum.

\subsubsection*{Angular resolution}

\begin{table}[H]
\centering\scriptsize
\caption{Q3b angular downsampling sweep. Closure error in metres.}
\label{tab:q3b_ang}
\resizebox{\columnwidth}{!}{%
\begin{tabular}{lccc}
\toprule
Setting & Basement\_1 & Floor7\_Hallway & Outdoor\_1 \\
\midrule
All beams      & \textbf{0.24} & 1.08 & 7.15 \\
Every 2nd beam & 0.68          & \textbf{0.71} & 3.95 \\
Every 3rd beam & 0.26          & 1.21 & 4.76 \\
\bottomrule
\end{tabular}%
}
\end{table}

\begin{figure}[H]
  \centering
  \includegraphics[width=\columnwidth]{plots/q3b_grids_basement_1_angular_resolution.png}
  \caption{Basement\_1 grids at full, 2nd, 3rd beam.}
  \label{fig:q3b_ang}
\end{figure}

Per-point normals come from PCA on $k = 5$ nearest neighbours. With
every beam those neighbours sit within a narrow angular window and
produce a sharp normal estimate. With every 3rd beam the same 5
neighbours span a much wider arc, the local surface is not well
represented, and the normal direction becomes noisier. Point-to-plane
ICP is highly sensitive to normal noise: a 5 degree normal error on a
5\,m wall translates to roughly 0.4\,m of along-wall slip per ICP
step. Floor7\_Hallway is the non-monotonic case, best at every 2nd
beam (0.71\,m) rather than at full resolution, because the tightly
spaced corridor walls cause ICP to over-fit micro-features at full
resolution and the every-2nd-beam setting smooths out the micro-noise
without losing the geometric constraint. Basement\_1, with cleaner
walls, prefers full resolution. There is no universal best setting;
the choice depends on environment scale.

\subsubsection*{Voxel grid downsampling}

\begin{table}[H]
\centering\scriptsize
\caption{Q3b voxel-size sweep. Closure error in metres.}
\label{tab:q3b_voxel}
\resizebox{\columnwidth}{!}{%
\begin{tabular}{lccc}
\toprule
Voxel & Basement\_1 & Floor7\_Hallway & Outdoor\_1 \\
\midrule
None    & \textbf{0.24} & 1.08 & 7.15 \\
0.05\,m & 0.41          & \textbf{0.14} & 12.4 \\
0.10\,m & 0.29          & 3.04 & 6.22 \\
0.20\,m & 2.60          & 2.94 & 6.02 \\
\bottomrule
\end{tabular}%
}
\end{table}

\begin{figure}[H]
  \centering
  \includegraphics[width=\columnwidth]{plots/q3b_grids_basement_1_voxel_downsampling.png}
  \caption{Basement\_1 grids at four voxel sizes.}
  \label{fig:q3b_voxel}
\end{figure}

Voxel downsampling replaces each $v \times v$ spatial cell with one
centroid point. At $v = 0.05$\,m the cell is smaller than typical
RPLidar range noise ($\pm 5$ to $\pm 20$\,mm), so the centroid acts
as an effective denoiser. At $v \ge 0.10$\,m the cell starts to span
real geometric features; the PCA normal computed on those merged
centroids averages multiple wall orientations into a single tilted
vector, ICP drives each source point along the wrong direction, and
point-to-plane residuals collapse toward point-to-point residuals
(which lose the tangential constraint that keeps the trajectory
aligned with walls). Basement\_1 is best with no voxel filter
(0.24\,m closure) because the high-density raw returns help corner
alignment, while Floor7\_Hallway is best at 0.05\,m (0.14\,m closure)
because the corridor walls are flat and the denoising helps.
Outdoor\_1 has no good voxel setting because the underlying geometry
is too sparse for any setting to recover.

\subsubsection*{Scan rate}

\begin{table}[H]
\centering\scriptsize
\caption{Q3b scan-rate sweep. Closure error in metres.}
\label{tab:q3b_rate}
\resizebox{\columnwidth}{!}{%
\begin{tabular}{lccc}
\toprule
Setting & Basement\_1 & Floor7\_Hallway & Outdoor\_1 \\
\midrule
All scans (10\,Hz)        & \textbf{0.24} & 1.08 & 7.15 \\
50\,\% scans (skip every other) & 0.23 & 2.94 & 11.5 \\
33\,\% scans (skip 2 of 3)      & 0.26 & \textbf{0.73} & 10.9 \\
\bottomrule
\end{tabular}%
}
\end{table}

\begin{figure}[H]
  \centering
  \includegraphics[width=\columnwidth]{plots/q3b_grids_basement_1_scan_rate.png}
  \caption{Basement\_1 grids at three scan rates.}
  \label{fig:q3b_rate}
\end{figure}

At 0.4\,m/s walking speed and 10\,Hz scan rate the sensor moves about
4\,cm per scan. Skipping every other scan brings this to 8\,cm and
skipping 2 of 3 brings it to 12\,cm; ICP's convergence basin is
roughly 0.5\,m, so all three settings stay well inside it on
Basement\_1 and the sweep is flat there. Floor7\_Hallway is the
non-monotonic case again, benefiting from 33\,\% scan rate (0.73\,m
versus 1.08\,m at full rate) because the longer per-scan motion gives
more parallax between consecutive scans and helps ICP estimate
translation along the corridor axis (the aperture problem, Lab 08
Activity 3B). On Outdoor\_1 fewer scans means fewer chances to escape
the divergence rejection gate, and closure error grows.

\begin{figure*}[!ht]
  \centering
  \includegraphics[width=\textwidth]{plots/q3b_basement_1.png}
  \caption{Q3b full parameter analysis --- Basement\_1. Row~1: trajectory
    overlays per variant. Row~2: closure error. Row~3: map point density
    (lower = denser). Row~4: per-parameter observations and explanation.}
  \label{fig:q3b_summary}
\end{figure*}

\begin{figure*}[!ht]
  \centering
  \includegraphics[width=\textwidth]{plots/q3b_floor7_hallway.png}
  \caption{Q3b full parameter analysis --- Floor7\_Hallway. Same layout.
    Note the non-monotonic angular result (every~2nd beam best due to
    corridor micro-feature over-fitting) and the 33\,\% scan-rate
    benefit from aperture-problem parallax.}
  \label{fig:q3b_summary_floor7}
\end{figure*}

\begin{figure*}[!ht]
  \centering
  \includegraphics[width=\textwidth]{plots/q3b_outdoor_1.png}
  \caption{Q3b full parameter analysis --- Outdoor\_1. Trajectory freezes
    under every parameter setting; failure is geometric (open-courtyard
    aperture problem) not a tuning issue.}
  \label{fig:q3b_summary_outdoor}
\end{figure*}

The single most important parameter is the maximum range. Both
indoor sequences improve by a factor of 12 when the range is set to
the sensor maximum, which is a rank-deficiency argument from
point-to-plane ICP rather than a tuning preference. Outdoor\_1
fails under every parameter setting, confirming that the constraint
is geometric (open-space aperture problem) rather than a tuning
issue.

\subsection{Q3c: Loop closure detection}
\label{sec:q3c}

A candidate pair $(i, j)$ of keyframes passes through four gates in
order, each removing a different class of false positive. The first
gate is temporal separation $j - i \ge 10$ keyframes, which stops the
detector matching a keyframe to itself or to its very close
neighbours. The second is an adaptive arc gate
$\text{arc}(j) - \text{arc}(i) \ge \min(8\,\text{m},\,
0.4 \cdot \text{arc}_{\text{total}})$, which stops the detector
firing on a sliding window that has barely moved; the adaptive form
lets short sequences (where 8\,m would not qualify any pair) still
fire on long-enough revisits. The third is pose distance
$\|t_i - t_j\| \le 2.0$\,m, which stops the detector from trying to
align scans that are far apart in the current trajectory estimate.
The fourth is an ICP overlap score threshold of 0.70: after aligning scan $j$ into scan $i$'s frame, at least 70\,\% of scan $j$'s points must have a nearest neighbour within 0.4\,m in scan $i$. A genuine revisit registers cleanly and most points line up, so the overlap score is close to one; a false candidate produces poor alignment and fails this check.

A single physical revisit usually generates many candidate $(i, j)$
pairs because the local map overlaps heavily between adjacent
keyframes. We collapse each cluster to one factor by greedily keeping
the highest-scoring candidate and dropping any other within 15
keyframes in both $i$ and $j$, which stops the pose graph being
over-constrained by 10 to 20 near-duplicates of the same physical
loop.

\begin{lstlisting}[language=Python]
# loop detection (excerpt)
effective_arc = max(
    2.0, min(min_loop_arc_m, arc[-1]*0.4))
for j in range(min_separation, n):
    for i in range(0, j - min_separation + 1):
        if arc[j] - arc[i] < effective_arc:
            continue
        if dist_ij(i, j) > pose_dist_thresh:
            continue
        T_rel, loop_H, _ = icp_scan_to_map(
            src_local, ref, ref_nrm, init,
            return_hessian=True)
        score = (dists < icp_corr_thresh).mean()
        if score >= icp_score_thresh:
            raw.append((score, i, j,
                        T_rel, dist_ij, loop_H))
\end{lstlisting}

\begin{table}[H]
\centering\scriptsize
\caption{Q3c loop closures detected per sequence with the four-gate
filter.}
\label{tab:q3c}
\resizebox{\columnwidth}{!}{%
\begin{tabular}{lccc}
\toprule
Sequence & Candidates evaluated & Loops accepted \\
\midrule
Basement\_1     & 2933 & 10 \\
Floor7\_Hallway & 1991 &  8 \\
Outdoor\_1      &    8 &  0 \\
\bottomrule
\end{tabular}%
}
\end{table}

\begin{figure}[H]
  \centering
  \includegraphics[width=\columnwidth]{plots/q3c_basement_1.png}
  \caption{Q3c Basement\_1 loop arcs and score histogram.}
  \label{fig:q3c_basement}
\end{figure}

\begin{figure}[H]
  \centering
  \includegraphics[width=0.49\columnwidth]{plots/q3c_floor7_hallway.png}
  \hfill
  \includegraphics[width=0.49\columnwidth]{plots/q3c_outdoor_1.png}
  \caption{Q3c Floor7\_Hallway and Outdoor\_1.}
  \label{fig:q3c_others}
\end{figure}

The right-hand histograms in Figures~\ref{fig:q3c_basement} and
\ref{fig:q3c_others} are the empirical evidence for the 0.70
threshold. False matches (sliding-window pairs that share local
geometry but are not true revisits) cluster between scores 0.45 and
0.60, while true revisits cluster between 0.75 and 0.95. The clean
gap around 0.70 holds across both indoor sequences, so setting the
threshold at 0.70 keeps every confirmed revisit while rejecting
every false positive observed in our runs.

Outdoor\_1 detects no loops because the trajectory froze near scan
400. After that, every keyframe sits at almost the same 2D position,
so the arc gate rejects almost all $(i, j)$ pairs (only 8 candidates
pass gates 1 to 3) and none of those 8 candidates exceed the score
threshold. A frozen trajectory has no loop information to add, so the
detector correctly invents none.

\subsection{Q3d: Factor graph optimisation}
\label{sec:q3d}

Each keyframe $k$ becomes a Pose2 variable
$x_k = (t_x, t_y, \theta)_k$ in GTSAM. The first pose is anchored
with a tight Gaussian prior ($\sigma = 10^{-4}$ on every axis) to
fix the gauge. Two kinds of between-factors are added: odometry
factors connect every consecutive pair $(k, k+1)$ with measurement
$z_{k,k+1}$ from raw ICP and information matrix obtained as the
average of the two endpoint ICP Hessians, and loop closure factors
connect each accepted pair $(i, j)$ from Q3c with measurement
$z_{ij}$ from the loop-detection ICP and information matrix from the
loop ICP Hessian. The cost is the standard Mahalanobis sum of
pose-graph residuals, minimised by GTSAM's built-in
Levenberg-Marquardt solver. The ICP Hessian
is parameterised $(\theta, t_x, t_y)$ but GTSAM Pose2 expects
$(t_x, t_y, \theta)$, so we permute the matrix before passing it in.

\begin{lstlisting}[language=Python]
# pose graph (excerpt)
graph.add(PriorFactorPose2(
    0, Pose2(x0, y0, th0),
    noiseModel.Diagonal.Sigmas([1e-4]*3)))
for i, j, z, info in factors:
    if info.ndim == 2:
        noise = noiseModel.Gaussian.Information(
            info)
    else:
        noise = noiseModel.Diagonal.Sigmas(info)
    graph.add(BetweenFactorPose2(
        i, j, Pose2(z[0], z[1], z[2]), noise))
params = gtsam.LevenbergMarquardtParams()
params.setMaxIterations(200)
optimizer = gtsam.LevenbergMarquardtOptimizer(
    graph, initial, params)
result = optimizer.optimize()
\end{lstlisting}

\begin{table}[H]
\centering\footnotesize
\caption{Q3d closure error before and after PGO.}
\label{tab:q3d}
\resizebox{\columnwidth}{!}{%
\begin{tabular}{lcccc}
\toprule
Sequence & Loops & Cost $i \to f$ & Closure $i$ & Closure $f$ \\
\midrule
Floor7\_Hallway & 8  & 69762 $\to$ 2752 & 0.166\,m & 0.221\,m \\
Basement\_1     & 10 & 61253 $\to$ 7349 & 0.170\,m & 0.204\,m \\
Outdoor\_1      & 0  & 0 $\to$ 0        & 8.084\,m & 8.084\,m \\
\bottomrule
\end{tabular}%
}
\end{table}

\begin{figure}[H]
  \centering
  \includegraphics[width=\columnwidth]{plots/q3d_floor7_hallway.png}
  \caption{Q3d Floor7\_Hallway before and after PGO.}
  \label{fig:q3d_floor7}
\end{figure}

\begin{figure}[H]
  \centering
  \includegraphics[width=\columnwidth]{plots/q3d_basement_1.png}
  \caption{Q3d Basement\_1 before and after PGO.}
  \label{fig:q3d_basement}
\end{figure}

\begin{figure}[H]
  \centering
  \includegraphics[width=\columnwidth]{plots/q3d_grid_basement_1.png}
  \caption{Q3d Basement\_1 grid before and after PGO.}
  \label{fig:q3d_grid}
\end{figure}

On both Basement\_1 and Floor7\_Hallway GTSAM reduces the total
graph cost by a factor of 8 to 25, but the results are degraded in
two ways that must be examined separately.

\textbf{Trajectory closure error.} The closure-error bars in
Figures~\ref{fig:q3d_floor7} and \ref{fig:q3d_basement} show that
the post-PGO closure error is worse than before PGO on every
sequence (Floor7: 0.166\,m $\to$ 0.221\,m; Basement\_1: 0.170\,m
$\to$ 0.204\,m). The optimised trajectory satisfies the loop
constraints better in a global Mahalanobis sense (total cost is low),
but the start-to-end Euclidean distance has grown because closure
error measures only one implied edge that the optimiser never
explicitly constrains.

\textbf{Occupancy grid quality.} The right panel of
Figure~\ref{fig:q3d_grid} shows a more serious consequence: after
PGO the Basement\_1 occupancy grid is visually destroyed. The clean
rectangular room walls visible before optimisation are scattered
across the full map extent; the grid is no longer recognisable as a
room. This happens because GTSAM redistributes large pose corrections
across the odometry chain, causing every previously registered scan
to be reprojected from a different pose. With inconsistent
registrations the ray-cast log-odds updates cancel rather than
reinforce, fragmenting the walls. A well-tuned PGO (with a looser
loop-edge information matrix or a robust kernel) would spread the
correction more smoothly and preserve wall coherence.

\textbf{Root cause.} The GTSAM solver emits the warning
\texttt{"Levenberg-Marquardt giving up because cannot decrease error
with maximum lambda"} on both sequences, meaning the loop and
odometry constraints are in tension and the solver cannot find a
descent direction that satisfies both simultaneously. The per-loop
ICP Hessian is over-confident (too tight a covariance) and forces
large abrupt pose jumps. Reducing the per-loop information weight
by a constant multiplier, scaling by ICP score so weak loops
contribute less, adding a Cauchy or Huber robust kernel on the loop
edges, or using chordal initialisation would each let the optimiser
trust the odometry chain more; we document the failure mode rather
than re-tune for the submitted figures, since the diagnostic itself
demonstrates understanding of pose-graph mechanics.

Outdoor\_1 is a degenerate case. With zero detected loops the pose
graph contains only odometry edges, the cost is identically zero
(every edge agrees with itself), and PGO returns the input unchanged.
This is the expected behaviour with no loop information: drift cannot
be corrected without a global anchor.

\subsection{Q3e and Q3f: video and presentation}

A single video file (\texttt{COMP0222\_CW2\_GRP\_32.mp4}) accompanies
the report. The LiDAR portion is the second half (last 30 seconds)
and shows, for each of the three sequences side by side, the
trajectory drawing itself in white over the accumulating colour point
cloud, with a green start marker and a red current marker. The map
builds up as the robot walks each loop, so the viewer can see the
map's consistency hold through the second loop on Basement\_1 and
Floor7\_Hallway, and see where Outdoor\_1 freezes. The Q3 video is
concatenated with the Q2 visual-SLAM video using \texttt{ffmpeg
concat} so that one merged file covers both questions. The five
minute oral presentation centres on the same material: data
collection, parameter analysis with Basement\_1 as the cleanest case,
loop detection and the threshold histogram, and the honest discussion
of the PGO degradation on the submitted sequences with the proposed
fix.

%% ===========================================================================
\section{Conclusion}
\label{sec:conclusion}

ORB-SLAM2 monocular on KITTI~07 reaches an ATE of 4.39\,m, which is
0.63\,\% of the 695\,m path and a competitive baseline for a
monocular run on this sequence. The two ablations on the same
benchmark probe different parts of the system. Disabling outlier
rejection on TUM blows up ATE by 1380\,\%, because local bundle
adjustment trusts every measurement and a single bad match shifts
the pose estimate immediately, so the next frame inherits the shift
and drift compounds frame by frame. The loop closure ablation on
KITTI adds 310\,\% to ATE while RPE actually drops, which is direct
evidence of the textbook ATE versus RPE distinction: per-step motion
is fine, but without the global anchor provided by Sim(3) loop
closure the heading drift accumulated over 695\,m never gets
cancelled. Cutting the feature budget below the default never helps,
and the one apparent win on TUM (feat=500, ATE 0.015\,m) is a
tracking failure dressed as a metric improvement, because only 45\,\%
of frames are tracked and ATE is measured on tracked frames only.

For the custom monocular sequences in Q2 we report inter-method
agreement between ORB-SLAM2 and COLMAP, because no external ground
truth exists. Basement\_1 is the clean indoor case at 132\,mm
translation disagreement after Umeyama scale alignment, while
Outdoor\_1 disagrees by 1.84\,m and 169 degrees because monocular
scale is poorly observable in open scenes and the two pipelines
settle on different essential-matrix branches.

For 2D LiDAR, the single most important parameter is the maximum
range. Truncating the RPLidar A2M12 from 12\,m to 2\,m inflates the
Basement\_1 closure from 0.19\,m to 2.38\,m, a factor of 12, because
ICP loses the opposite-wall geometry it needs to constrain
translation. The pose graph optimisation in Q3d is the most
informative finding of the project: GTSAM reduces the total
Mahalanobis cost by a factor of 8 to 25 with per-loop ICP-Hessian
information matrices, but the closure error grows slightly because
the loop edges and odometry edges end up in tension and the
Levenberg-Marquardt solver hits its lambda ceiling. The fix is to
soften the loop weighting (a Cauchy or Huber kernel on loop edges, or
a constant down-weight on the loop information matrix) so that the
optimiser trusts the odometry chain more and the closure-improving
solution becomes reachable.

\end{document}
